% Generated draft: requires mathematical and source review.
\documentclass[11pt]{article}
\usepackage[margin=25mm]{geometry}
\usepackage{fontspec}
\setmainfont{Latin Modern Roman}
\usepackage{amsmath,amssymb,mathtools,unicode-math}
\setmathfont{Latin Modern Math}
\usepackage{enumitem,needspace,xurl,hyperref}
\setlength{\parindent}{0pt}
\setlength{\parskip}{5pt}
\setlength{\emergencystretch}{3em}
\newcommand{\N}{\mathbb{N}}
\newcommand{\Z}{\mathbb{Z}}
\newcommand{\Q}{\mathbb{Q}}
\newcommand{\R}{\mathbb{R}}
\newcommand{\C}{\mathbb{C}}
\newcommand{\F}{\mathbb{F}}
\newcommand{\A}{\mathbb{A}}
\newcommand{\PP}{\mathbb{P}}
\newcommand{\E}{\mathbb{E}}
\newcommand{\eps}{\varepsilon}
\newcommand{\dd}{\,\mathrm{d}}
\begin{document}
\section*{UnsolvedMath ID: 30004889}
\textbf{Source problem code:} catalogue entry
\subsection{Definitions and mathematical statement}
The supplied registry proposes the following unitarity question. Its formulation depends on an unspecified meaning of ``unipotent ideal associated with a nilpotent orbit'' and on an unspecified orbit convention; these qualifications remain unresolved.

Let \(G\) be a real reductive Lie group, let \(\mathfrak g_{\R}=\operatorname{Lie}(G)\), put \(\mathfrak g=\mathfrak g_{\R}\otimes_{\R}\C\), and let \(K\) be a maximal compact subgroup. Write \(U(\mathfrak g)\) for the universal enveloping algebra. A Harish--Chandra module is a finitely generated \(U(\mathfrak g)\)-module with a compatible, locally finite \(K\)-action, finite multiplicities of irreducible \(K\)-types, and agreement of the differentiated \(K\)-action with the restricted Lie algebra action. The precise class of real reductive groups intended by the record is not specified.

Let \(\mathcal O\) be a nilpotent orbit in the intended, unspecified orbit setting, and let \(I\subseteq U(\mathfrak g)\) be an ideal designated as a unipotent ideal associated with \(\mathcal O\). The record supplies neither a definition of this designation nor the association. Relative to such a choice, an irreducible Harish--Chandra module \(X\) is called unipotent when
\[
I X=0,
\qquad\text{equivalently}\qquad
I\subseteq\operatorname{Ann}_{U(\mathfrak g)}(X).
\]
This condition does not require equality of the two ideals.

For an orbit-closure interpretation, write
\[
\partial\mathcal O=\overline{\mathcal O}\setminus\mathcal O,
\qquad
c(\mathcal O)=\dim\mathcal O-\dim\partial\mathcal O.
\]
Here the closure topology and the choice of real or complex dimension must be specified consistently with the intended orbit setting. The expression \(\operatorname{codim}(\partial\mathcal O,\mathcal O)\) in the record is read as this boundary codimension in the orbit closure, since the boundary is not a subset of the orbit itself. The record gives no convention when the boundary is empty.

With these unresolved choices fixed, the proposed universal question is: for every such \(G\), every associated pair \((\mathcal O,I)\), and every nonzero irreducible Harish--Chandra module \(X\), does
\[
I X=0\quad\text{and}\quad c(\mathcal O)\geq 4
\quad\Longrightarrow\quad
X\text{ is unitary}
\]
hold? In the algebraic unitarity interpretation, the conclusion means that \(X\) admits a positive definite Hermitian form \(h\) satisfying
\[
h(kv,kw)=h(v,w),\qquad
h(Av,w)+h(v,Aw)=0
\]
for all \(k\in K\), \(A\in\mathfrak g_{\R}\), and \(v,w\in X\). The record does not specify whether its intended conclusion additionally requires a Hilbert space globalization as a unitary representation of the stated group \(G\).

The supplied original excerpt concerns unipotent \emph{bimodules}, asserts a special form for their underlying \(G\)-representations in connection with a conjecture of Vogan labelled [V3], and conjectures that an unidentified preceding problem disappears under the boundary-codimension condition. It does not explicitly state the unitarity implication above. Thus the displayed question records the supplied cleaned formulation, but its equivalence to the original intended conjecture cannot be established from this record.
\subsection{Short English statement}
The cleaned record asks whether every irreducible unipotent Harish--Chandra module is unitary when the associated nilpotent orbit has boundary codimension at least four. The meaning of the associated unipotent ideal, the orbit and codimension conventions, and the connection to the original bimodule excerpt require clarification. This is a statement-only draft requiring human mathematical review.
\subsection{Sources}
\begin{itemize}
\item Local registry: \texttt{data/unsolvedmath/problems.json}, permanent UnsolvedMath ID \(30004889\), source problem code \texttt{catalogue entry}, title ``Unitarity from Nilpotent-Orbit Boundary Codimension.'' The supplied collection is ``Oberwolfach Reports Open Problems.'' No individual proposer is identified; the proposed year is recorded as 2021.
\item Supplied report citation: \emph{Enveloping Algebras and Geometric Representation Theory} (2022), Oberwolfach Reports, pp.\ 2827--2891, DOI \texttt{10.4171/owr/2021/52}. \url{https://doi.org/10.4171/owr/2021/52}. The extracted original passage carries page number 2851. The report title, publication year, DOI, and page references are preserved as supplied.
\item Supplied literature source labelled ``Mason-Brown'': \url{https://arxiv.org/abs/1805.12038}.
\item Supplied literature source labelled ``Losev--Mason-Brown--Matvieievskyi'': \url{https://arxiv.org/abs/2108.03453}.
\item Supplied literature source labelled ``OWR report'': \url{https://doi.org/10.4171/owr/2021/52}.
\item Literature assessment dated exactly \texttt{2026-08-22}: the supplied curation labels the problem ``partially solved'' and states that results prove Vogan-type structure statements under codimension at least four and establish broad classical-group cases, while the exact unitarity conjecture for all real reductive groups remains unsettled. This is supplied dated curation, not independently verified here; the record itself requests independent verification of time-sensitive status and cited claims.
\item The registry labels its statement \texttt{corrected\_verified} and reports that the official Oberwolfach report was checked directly and that the clean statement preserves the complete target while restoring definitions and removing extraction artifacts or source-dependent wording. This is a supplied verification claim, not a verification performed for this draft. The original excerpt attributes a conjecture to Vogan using the unresolved reference label [V3]; no corresponding bibliographic entry is supplied.
\end{itemize}
\end{document}
